% !TEX root = ../main.tex
\chapter{THIẾT KẾ VÀ HIỆN THỰC}

Dựa trên kết quả phân tích ở Chương~3, chương này trình bày cách hệ thống được hiện thực cụ thể: từ thiết kế cơ sở dữ liệu, kiến trúc phần mềm đến giao diện người dùng và kết quả kiểm thử. Toàn bộ hệ thống được xây dựng trên nền tảng NestJS~\cite{nestjs-docs}, Next.js~\cite{nextjs-docs} và React Native~\cite{reactnative-docs}.

\section{Sơ đồ lớp}

Để tăng tính rõ ràng, sơ đồ lớp được tách thành hai nhóm: nhóm \textbf{Core} phụ trách tài khoản và vận hành nền tảng; nhóm \textbf{IELTS} phụ trách toàn bộ nội dung và tiến độ học tập~\cite{fowler2002patterns}.

Nhóm Core tập trung vào quản lý tài khoản và xác thực. Lớp User là trung tâm, lưu các thông tin cơ bản như email, mật khẩu (mã hoá bằng bcrypt), họ tên, avatar và vai trò (STUDENT / INSTRUCTOR / ADMIN). Người dùng đăng nhập bằng Google OAuth sẽ không có trường mật khẩu. Danh tính được xác thực qua JWT. Lớp Subscription liên kết 1-1 với User để quản lý gói cước hiện tại (FREE / PREMIUM / PRO) và ngày hết hạn.

Nhóm IELTS chứa toàn bộ nội dung học tập, tổ chức theo từng giai đoạn trong lộ trình~\cite{fowler2002patterns}. Ở giai đoạn Basic, bốn kỹ năng được tách thành bốn lớp exercise riêng biệt (Listening, Reading, Writing, Speaking) kèm bảng tiến độ. Giai đoạn Advanced bổ sung thêm lớp Session để lưu lại từng lần luyện tập cùng điểm số phản hồi từ AI. Giai đoạn Intensive dùng lớp IeltsIntensiveSession để quản lý phiên thi thử với vòng đời trạng thái từ khi bắt đầu cho đến khi có kết quả chấm điểm. Ngoài ra, phân hệ Vocab Lab dùng lớp Flashcard với thuật toán nhắc lại FSRS (qua thư viện ts-fsrs) để giúp người học ghi nhớ từ vựng theo lịch tối ưu.

\begin{figure}[H]
    \centering
    \includegraphics[width=\linewidth]{4.1. Class Diagram/2-class-core.png}
    \caption{Sơ đồ lớp nhóm Core}
    \label{fig:class-user}
\end{figure}

\begin{figure}[H]
    \centering
    \includegraphics[width=\linewidth]{4.1. Class Diagram/1-class-ielts.png}
    \caption{Sơ đồ lớp nhóm IELTS}
    \label{fig:class-learning}
\end{figure}

\section{Sơ đồ cơ sở dữ liệu}

Hệ thống dùng PostgreSQL~\cite{postgresql-docs} thông qua Prisma ORM~\cite{prisma-docs}. Cơ sở dữ liệu được lưu trữ trên \textbf{Supabase} — một dịch vụ PostgreSQL được quản lý sẵn, kết nối qua PgBouncer (port 6543) trong môi trường chạy thực và kết nối trực tiếp (port 5432) khi chạy migration. Toàn bộ schema được định nghĩa tập trung trong file schema.prisma với 66 model. Tương tự sơ đồ lớp, ERD cũng được tách thành hai nhóm để dễ đọc: nhóm \textbf{Core} và nhóm \textbf{IELTS}.

Nhóm Core quản lý người dùng và các nghiệp vụ liên quan. Bảng users lưu thông tin tài khoản, trạng thái hoạt động và định danh Google OAuth. Bên cạnh đó là các bảng subscriptions, pricing\_plans (gói cước), payments (thanh toán), notifications (thông báo) và student\_teacher\_links (liên kết học viên -- giáo viên). Sơ đồ ERD nhóm Core được trình bày trong hình~\ref{fig:erd-core}.

\begin{figure}[H]
    \centering
    \includegraphics[width=\linewidth]{4.2. ERD/2-erd-core.png}
    \caption{Sơ đồ ERD nhóm Core (PostgreSQL)}
    \label{fig:erd-core}
\end{figure}

Nhóm IELTS chứa toàn bộ dữ liệu học thuật, được tổ chức theo bốn giai đoạn của lộ trình. Mỗi giai đoạn có nhóm bảng riêng:

\begin{itemize}
    \item \textbf{Foundation:} Bộ ba nhóm bảng từ vựng nền, ngữ pháp và phát âm. Mỗi nhóm có cấu trúc phân cấp từ cuốn sách, chương, đơn vị bài tập cho đến bảng ghi nhận tiến độ của từng người học.
    \item \textbf{Basic:} Bốn bảng bài tập tương ứng với bốn kỹ năng Listening, Reading, Writing, Speaking, cùng bảng tiến độ để theo dõi kết quả.
    \item \textbf{Advanced:} Mỗi kỹ năng có thêm bảng Session lưu lại từng lần luyện tập kèm điểm phản hồi từ AI.
    \item \textbf{Intensive:} Ba bảng chính là exams (đề thi), exam\_sessions (phiên thi với vòng đời trạng thái) và results (kết quả chấm điểm AI sau khi xử lý).
\end{itemize}

Ngoài ra, phân hệ Vocab Lab dùng các bảng decks, flashcards và flashcard\_reviews để lưu trạng thái FSRS của từng thẻ. Các bảng Dictation, Shadowing và lộ trình học cũng được bổ sung để phục vụ các tính năng còn lại. Khoá ngoại được thiết lập với ON DELETE CASCADE ở toàn bộ bảng con nhằm giữ tính toàn vẹn dữ liệu khi xoá bản ghi cha. Sơ đồ ERD nhóm IELTS được trình bày trong hình~\ref{fig:erd-learning}.

\begin{figure}[H]
    \centering
    \includegraphics[width=\linewidth]{4.2. ERD/1-erd-ielts.png}
    \caption{Sơ đồ ERD nhóm IELTS (PostgreSQL)}
    \label{fig:erd-learning}
\end{figure}

\section{Sơ đồ kiến trúc phần mềm}

Hệ thống được xây dựng theo kiến trúc \textbf{Event-driven Modular Monolith}~\cite{newman2015microservices}, gồm các thành phần chính như sau:

\begin{itemize}
    \item \textbf{NestJS Backend Core:} Toàn bộ nghiệp vụ được tổ chức thành 17 module trong một ứng dụng NestJS 10 duy nhất. Các module giao tiếp nội bộ qua Dependency Injection, giúp dễ bảo trì mà không cần phân tán thành nhiều service riêng lẻ.
    \item \textbf{FastAPI AI Worker:} Một microservice Python độc lập, chịu trách nhiệm cho toàn bộ tác vụ AI. Worker này nhận công việc từ NestJS qua hàng đợi RabbitMQ (CloudAMQP), gọi Gemini API để chấm điểm, rồi ghi kết quả trực tiếp vào PostgreSQL qua psycopg2 bằng câu truy vấn SQL thô.
    \item \textbf{Tầng dữ liệu:} PostgreSQL trên Supabase (kết nối qua PgBouncer), Redis trên Upstash (cache và phiên làm việc), Cloudinary lưu trữ ảnh và file audio ghi âm của người dùng, và Google Cloud Storage phục vụ lưu trữ media tĩnh cho AI Worker khi chấm điểm.
    \item \textbf{Client:} Ứng dụng web Next.js và ứng dụng mobile React Native, đều kết nối với Backend Core qua HTTPS thông qua nginx.
\end{itemize}

\begin{figure}[H]
    \centering
    \includegraphics[width=\linewidth]{4.3. Architecture Diagram/architecture-diagram.png}
    \caption{Sơ đồ kiến trúc phần mềm hệ thống IELTS Master}
    \label{fig:architecture}
\end{figure}

\section{Sơ đồ luồng màn hình}

\subsection{Sơ đồ luồng màn hình website}

Trên nền tảng Web, người dùng bắt đầu từ trang chủ Landing Page, đăng nhập hoặc đăng ký tài khoản, rồi vào Dashboard chính. Tại đây, người học theo lộ trình từ Foundation (xây dựng nền tảng từ vựng, ngữ pháp, phát âm), qua Basic (làm quen bốn kỹ năng cơ bản), Advanced (luyện chuyên sâu với phản hồi từ AI), đến Intensive (thi thử sát thực tế). Ngoài ra còn có các công cụ bổ trợ như Vocab Lab (sổ tay từ vựng cá nhân), Dictation, Shadowing, Community, Roadmap và trang Thống kê để theo dõi tiến độ. Người quản lý có nhánh riêng vào Admin Dashboard với các chức năng quản lý nội dung, đề thi và tài khoản học viên.

\begin{figure}[H]
    \centering
    \includegraphics[width=\linewidth]{4.4. Flow Chart/1-flowchart-web.png}
    \caption{Sơ đồ luồng màn hình website (Next.js)}
    \label{fig:flow-web}
\end{figure}

\subsection{Sơ đồ luồng màn hình mobile}

Trên nền tảng Mobile (React Native), cấu trúc điều hướng sử dụng Tab Navigation ở tầng dưới cùng với 5 tab hiển thị: \textbf{Home} (Trang chủ), \textbf{Explore} (Khám phá), \textbf{IELTS} (trung tâm học tập), \textbf{Community} (Cộng đồng) và \textbf{Profile} (Hồ sơ). Các màn hình chuyên biệt như Vocab Lab, Shadowing, Pronunciation, Vocabulary, Grammar và More được truy cập qua điều hướng nội bộ từ tab IELTS và Home, không hiển thị trực tiếp trên thanh tab. Các màn hình con được quản lý bằng Stack Navigator tích hợp trong Expo Router.

\begin{figure}[H]
    \centering
    \includegraphics[width=\linewidth]{4.4. Flow Chart/2-flowchart-mobile.png}
    \caption{Sơ đồ luồng màn hình mobile (React Native)}
    \label{fig:flow-mobile}
\end{figure}

\section{Giao diện chương trình}

Phần này trình bày 13 giao diện tiêu biểu của hệ thống trên nền tảng Web, sắp xếp theo hành trình người dùng từ khi lần đầu truy cập cho đến khi sử dụng các tính năng học tập chuyên sâu.

\subsection{Giao diện trang chủ (Landing Page)}

Trang chủ là màn hình đầu tiên người dùng nhìn thấy khi truy cập hệ thống. Giao diện truyền tải thông điệp \textit{``Master English, Ace IELTS, Smarter with AI''} qua bố cục hiện đại với hình ảnh minh hoạ sinh động và hiệu ứng hoạt hình nhẹ. Nút ``START LEARNING'' đặt nổi bật ở trung tâm để khuyến khích người dùng bắt đầu ngay. Khi nhấn vào đó, hệ thống đưa người dùng thẳng vào phân hệ học tập chính. Giao diện trang chủ được trình bày trong hình~\ref{fig:ui-landing}.

\begin{figure}[H]
    \centering
    
\includegraphics[width=\linewidth]{4.5. Program Interface/images/01-homepage.png}
    \caption{Giao diện trang chủ (Landing Page)}
    \label{fig:ui-landing}
\end{figure}

\subsection{Giao diện Dashboard người học}

Trang IELTS Dashboard đóng vai trò là trung tâm điều khiển (hub) chính, nơi học viên theo dõi và quản lý toàn bộ quá trình luyện thi IELTS của mình. Mục đích cốt lõi của trang là cung cấp một Lộ trình học tập (Learning Roadmap) được cấu trúc rõ ràng qua 4 giai đoạn tiến triển: Foundation (xây dựng nền tảng từ vựng, ngữ pháp, phát âm), Basic (nắm bắt kiến thức cơ bản của 4 kỹ năng), Advanced (luyện tập chuyên sâu từng dạng câu hỏi với AI), và Intensive (thi thử sát với thực tế). Ngoài việc hiển thị thanh tiến độ và số liệu tổng quan (Quick Stats), Dashboard còn đóng vai trò là điểm truy cập cho các công cụ bổ trợ (Tools \& Features) quan trọng như tạo lộ trình học cá nhân hóa, Công cụ tính điểm Band Score, Kết nối theo dõi tiến độ giữa Học viên và Giáo viên, cùng với hệ thống Thống kê chi tiết giúp đánh giá tổng quan năng lực và nhận diện những kỹ năng cần cải thiện. Giao diện Dashboard được trình bày trong hình~\ref{fig:ui-dashboard}.

\begin{figure}[H]
    \centering
    \includegraphics[width=\linewidth]{4.5. Program Interface/images/02-ielts-dashboard.png}
    \caption{Giao diện Dashboard người học}
    \label{fig:ui-dashboard}
\end{figure}

\subsection{Giao diện Lộ trình học IELTS (Roadmap)}

Trang IELTS Roadmap là công cụ hỗ trợ xây dựng lộ trình học tập cá nhân hóa. Dựa trên kết quả Placement Test, mục tiêu Band Score và lịch học của người dùng, hệ thống tự động tạo lộ trình ôn luyện chi tiết, ưu tiên các kỹ năng còn yếu và bỏ qua những kỹ năng đã đạt ngưỡng thành thạo. Người học theo dõi tiến độ hoàn thành từng bước trực tiếp trên trang này; trạng thái từng bước được cập nhật tự động khi bài tập tương ứng được hoàn thành. Giao diện Roadmap được trình bày trong hình~\ref{fig:ui-roadmap}.

\begin{figure}[H]
    \centering
    \includegraphics[width=\linewidth]{4.5. Program Interface/images/09-ielts-roadmap.png}
    \caption{Giao diện Lộ trình học IELTS cá nhân hoá (Roadmap)}
    \label{fig:ui-roadmap}
\end{figure}

\subsection{Giao diện thư viện Từ vựng (Foundation Vocabulary)}

Trang Vocabulary (Từ vựng) thuộc giai đoạn Foundation, đóng vai trò là thư viện sách từ vựng hỗ trợ học viên xây dựng vốn từ cơ bản. Mục đích chính của trang này là hiển thị danh sách các ``cuốn sách'' từ vựng theo chủ đề (Foundation Vocab Book), bao gồm hình ảnh minh họa, tên chủ đề và tổng số lượng từ vựng trong mỗi sách. Người dùng có thể dễ dàng duyệt qua các chủ đề và bấm nút ``START LEARNING'' để chuyển trực tiếp vào màn hình học chi tiết của từng chủ đề đó, tạo bước đệm vững chắc trước khi bước vào luyện thi IELTS chuyên sâu. Giao diện thư viện Vocabulary được trình bày trong hình~\ref{fig:ui-vocabulary}.

\begin{figure}[H]
    \centering
    \includegraphics[width=\linewidth]{4.5. Program Interface/images/03-vocabulary.png}
    \caption{Giao diện thư viện Từ vựng giai đoạn Foundation}
    \label{fig:ui-vocabulary}
\end{figure}

\subsection{Giao diện luyện Listening cơ bản (Basic Listening)}

Trang IELTS Basic Listening Exercise là giao diện làm bài tập thực hành kỹ năng Nghe ở giai đoạn cơ bản. Trải nghiệm người dùng ở đây tập trung vào việc nghe audio và tương tác trả lời câu hỏi theo các định dạng bài thi quen thuộc (như điền từ, trắc nghiệm). Sau khi nộp bài, hệ thống sẽ tự động chấm điểm và cung cấp phản hồi (đáp án đúng, transcript) ngay lập tức để học viên rút kinh nghiệm và củng cố kỹ năng nghe cho từng dạng bài cụ thể. Giao diện luyện Listening cơ bản được trình bày trong hình~\ref{fig:ui-listening-basic}.

\begin{figure}[H]
    \centering
    \includegraphics[width=\linewidth]{4.5. Program Interface/images/04-ielts-basic-listening-exercise.png}
    \caption{Giao diện luyện Listening giai đoạn Basic}
    \label{fig:ui-listening-basic}
\end{figure}

\subsection{Giao diện thi thử Reading (Intensive Reading)}

Trang IELTS Intensive Reading Test Taken cung cấp môi trường thi thử giả lập sát với bài thi Reading thực tế. Giao diện được chia đôi: một bên là các đoạn văn bản học thuật (Passages), bên còn lại là danh sách câu hỏi cần trả lời. Điểm nhấn của trang này là tính năng mô phỏng áp lực thời gian (đồng hồ đếm ngược) và khả năng tương tác trực tiếp với văn bản (như highlight), giúp người dùng làm quen với nhịp độ và không khí của một bài thi IELTS Reading thực sự. Giao diện thi thử Reading được trình bày trong hình~\ref{fig:ui-reading-intensive}.

\begin{figure}[H]
    \centering
    \includegraphics[width=\linewidth]{4.5. Program Interface/images/05-ielts-intensive-reading-test-taken.png}
    \caption{Giao diện thi thử Reading (Intensive Reading Test)}
    \label{fig:ui-reading-intensive}
\end{figure}

\subsection{Giao diện trang đáp án Listening (Intensive Listening Answer)}

Trang IELTS Intensive Listening Answer Page là giao diện hiển thị kết quả chi tiết sau khi học viên hoàn thành bài thi thử Listening. Mục đích chính là cung cấp báo cáo điểm số (tổng điểm và Band Score dự kiến), đồng thời liệt kê chi tiết từng câu hỏi để người dùng đối chiếu đáp án đã chọn với đáp án đúng. Trang này thường tích hợp sẵn Audio Transcript có bôi đậm đoạn chứa câu trả lời, giúp người học dễ dàng phân tích lỗi sai và cải thiện kỹ năng nghe hiểu. Giao diện trang đáp án Listening được trình bày trong hình~\ref{fig:ui-listening-answer}.

\begin{figure}[H]
    \centering
    \includegraphics[width=\linewidth]{4.5. Program Interface/images/06-ielts-intensive-listening-answer-page.png}
    \caption{Giao diện trang đáp án Listening sau bài thi thử}
    \label{fig:ui-listening-answer}
\end{figure}

\subsection{Giao diện thi thử Speaking (Intensive Speaking)}

Trang IELTS Intensive Speaking Test Taken mô phỏng lại một phòng thi Speaking thực tế với sự hỗ trợ của AI. Người học sẽ tương tác với một ``giám khảo AI'', lắng nghe các câu hỏi thuộc Part 1, Part 2, Part 3 và sử dụng microphone để thu âm câu trả lời của mình trong thời gian quy định. Giao diện được thiết kế tối giản để người dùng tập trung vào việc nói, đồng thời có các chỉ báo thu âm rõ ràng để đảm bảo quá trình thi diễn ra mượt mà trước khi hệ thống gửi file âm thanh đi chấm điểm. Giao diện thi thử Speaking được trình bày trong hình~\ref{fig:ui-speaking}.

\begin{figure}[H]
    \centering
    \includegraphics[width=\linewidth]{4.5. Program Interface/images/07-ielts-intensive-speaking-test-taken.png}
    \caption{Giao diện thi thử Speaking với giám khảo AI}
    \label{fig:ui-speaking}
\end{figure}

\subsection{Giao diện kết quả Writing với phản hồi AI (Intensive Writing Answer)}

Trang IELTS Intensive Writing Answer Page có chức năng hiển thị kết quả đánh giá bài thi Viết của học viên sau khi được chấm điểm bởi hệ thống AI. Giao diện cung cấp cái nhìn toàn diện về bài làm, bao gồm Band Score dự kiến, nhận xét tổng quan và phân tích chi tiết theo 4 tiêu chí chấm thi của IELTS (Task Achievement, Coherence, Lexical Resource, Grammatical Range). Đặc biệt, trang này liệt kê các lỗi ngôn ngữ cụ thể được phát hiện (kèm câu gốc, phiên bản đã sửa và giải thích ngắn), giúp người dùng nhận ra điểm yếu về ngữ pháp và từ vựng để cải thiện bài viết. Giao diện kết quả Writing được trình bày trong hình~\ref{fig:ui-writing-answer}.

\begin{figure}[H]
    \centering
    \includegraphics[width=\linewidth]{4.5. Program Interface/images/08-ielts-intensive-writing-answer-page.png}
    \caption{Giao diện kết quả Writing với phản hồi chi tiết của AI}
    \label{fig:ui-writing-answer}
\end{figure}

\subsection{Giao diện luyện nghe chép chính tả (Dictation)}

Trang Dictation Lessons (Luyện nghe chép chính tả) là một phân hệ luyện tập chuyên biệt giúp nâng cao sự nhạy bén của kỹ năng nghe. Tại đây, người dùng có thể nghe các đoạn audio ngắn và thực hành gõ lại toàn bộ nội dung (Dictation). Giao diện chú trọng công cụ đối chiếu độ chính xác giữa bài làm của người dùng với transcript gốc. Giao diện luyện Dictation được trình bày trong hình~\ref{fig:ui-dictation}.

\begin{figure}[H]
    \centering
    \includegraphics[width=\linewidth]{4.5. Program Interface/images/10-dictation-lessons.png}
    \caption{Giao diện luyện nghe chép chính tả (Dictation)}
    \label{fig:ui-dictation}
\end{figure}

\subsection{Giao diện sổ tay từ vựng cá nhân (Vocab Lab)}

Trang Vocab Lab đóng vai trò như một ``phòng thí nghiệm'' hoặc cuốn sổ tay từ vựng kỹ thuật số cá nhân của người học. Mục đích chính là lưu trữ, quản lý và ôn tập những từ vựng mà học viên đã lưu lại trong suốt quá trình học trên hệ thống. Trang này cung cấp các tính năng như tạo danh sách từ vựng tùy chỉnh, ôn tập qua flashcard và áp dụng thuật toán Lặp lại ngắt quãng (Spaced Repetition) để giúp người học ghi nhớ từ vựng một cách khoa học và lâu dài. Giao diện Vocab Lab được trình bày trong hình~\ref{fig:ui-vocablab}.

\begin{figure}[H]
    \centering
    \includegraphics[width=\linewidth]{4.5. Program Interface/images/11-vocab-lab.png}
    \caption{Giao diện sổ tay từ vựng cá nhân (Vocab Lab)}
    \label{fig:ui-vocablab}
\end{figure}

\subsection{Giao diện Cộng đồng (Community)}

Trang Community (Cộng đồng) là không gian kết nối xã hội dành cho người dùng trên toàn hệ thống. Tại đây, học viên có thể chia sẻ mẹo học tập, thảo luận về kinh nghiệm ôn thi, đăng bài chia sẻ thành tích band điểm, và tương tác với các thành viên khác. Tính năng này giúp tạo ra một môi trường học tập mang tính cộng đồng, tăng cường động lực và sự gắn kết của người dùng với hệ thống. Giao diện Cộng đồng được trình bày trong hình~\ref{fig:ui-community}.

\begin{figure}[H]
    \centering
    \includegraphics[width=\linewidth]{4.5. Program Interface/images/12-community.png}
    \caption{Giao diện Cộng đồng (Community)}
    \label{fig:ui-community}
\end{figure}

\subsection{Tính năng tra từ và lưu vào Vocab Lab bằng AI}

Tính năng Add to Vocab Lab with AI là một panel nổi xuất hiện khi người học bôi đen từ hoặc cụm từ trên các màn hình học hỗ trợ (Foundation Vocabulary, IELTS Basic, Dictation). Panel cung cấp hai tùy chọn: \textit{``Add to Vocab''} điền sẵn từ và ngữ cảnh vào form tạo thẻ trong tab Add của Vocab Lab để người học bổ sung nội dung còn lại và lưu; \textit{``Ask AI''} mở Chat AI để trợ lý tự động sinh đầy đủ nội dung cho thẻ theo Card Type đang dùng. Người học có thể xem lại và chỉnh sửa trước khi lưu mà không cần rời khỏi trang đang học. Giao diện tính năng được trình bày trong hình~\ref{fig:ui-add-vocablab-ai}.

\begin{figure}[H]
    \centering
    \includegraphics[width=\linewidth]{4.5. Program Interface/images/13-add-to-vocab-lab-with-ai.png}
    \caption{Pop-up tra từ và lưu vào Vocab Lab bằng AI}
    \label{fig:ui-add-vocablab-ai}
\end{figure}

\section{Kiểm thử hệ thống}

Tầng backend-core của hệ thống được kiểm thử tự động với hai mục tiêu: phát hiện sớm regression sau mỗi commit và làm tài liệu sống mô tả hành vi mong đợi của từng mô-đun. Toàn bộ test case được tổ chức theo tháp kiểm thử ba tầng (unit, integration và end-to-end), thực thi tự động trong pipeline CI và sinh báo cáo Markdown để đối chiếu trực tiếp với báo cáo khóa luận này.

\subsection{Chiến lược và công cụ kiểm thử}

Khung kiểm thử chính là Jest 29.7.0 kết hợp Supertest cho các yêu cầu HTTP, ts-jest cho mã TypeScript và jest-junit để xuất kết quả ở định dạng JUnit phục vụ pipeline CI. Tầng unit cô lập từng service bằng cách mock toàn bộ phụ thuộc (PrismaService, JwtService, Redis, RabbitMQ) thông qua Dependency Injection của NestJS, qua đó chỉ kiểm chứng logic nghiệp vụ thuần. Tầng integration mock các phụ thuộc ngoài nhưng giữ nguyên Controller, Pipe, Guard và Interceptor bằng cách khởi tạo INestApplication thực; nhờ đó test xác minh được toàn bộ vòng đời request và response, bao gồm cả ValidationPipe, JwtAuthGuard và RolesGuard.

Tầng end-to-end (E2E) chạy độc lập bằng lệnh npm run test:e2e với cấu hình riêng jest-e2e.json. Khác với hai tầng trên, E2E sử dụng cơ sở dữ liệu PostgreSQL thật trên schema test thay vì mock, qua đó kiểm chứng được cả ràng buộc cấp cơ sở dữ liệu (chỉ mục duy nhất trên email, khóa ngoại) và hành vi stateless của JWT. Bộ E2E hiện gồm 12 test case mô phỏng toàn bộ luồng xác thực: đăng ký, đăng nhập, truy cập tài nguyên bảo vệ và refresh token; kèm các kịch bản tiêu cực như token hết hạn, token bị giả mạo chữ ký và đăng ký trùng email.

Để duy trì chất lượng mã nguồn theo thời gian, dự án áp dụng ngưỡng độ phủ tối thiểu (coverage threshold) trên các service trọng yếu, được khai báo trong tệp jest.config.ts. Khi lệnh npm run test:cov chạy trong pipeline CI, Jest trả về exit code khác 0 nếu coverage giảm xuống dưới ngưỡng, buộc CI fail và chặn merge pull request. Các ngưỡng hiện hành được trình bày trong bảng~\ref{tab:coverage_thresholds}.

% 3 cột p{}, divisor=100: 50+25+25; -4pt buffer cho rounding
\setlength{\tbu}{\dimexpr(\linewidth - 6\tabcolsep - 4\arrayrulewidth - 4pt)/100\relax}
\begin{longtable}{|p{50\tbu}|p{25\tbu}|p{25\tbu}|}
    \caption{Ngưỡng độ phủ mã (coverage threshold) backend-core}
    \label{tab:coverage_thresholds} \\
    \hline
    \textbf{Mô-đun / tệp nguồn} & \textbf{Lines tối thiểu (\%)} & \textbf{Branches tối thiểu (\%)} \\ \hline
    \endfirsthead
    \multicolumn{3}{r}{\textit{(tiếp trang trước)}} \\
    \hline
    \textbf{Mô-đun / tệp nguồn} & \textbf{Lines tối thiểu (\%)} & \textbf{Branches tối thiểu (\%)} \\ \hline
    \endhead
    \hline \multicolumn{3}{r}{\textit{(tiếp trang sau)}} \\
    \endfoot
    \hline
    \endlastfoot

    auth.service.ts (xác thực và phân quyền) & 70 & 60 \\ \hline
    subscriptions.service.ts (lifecycle gói Premium và thanh toán VNPay) & 50 & 30 \\ \hline
    gamification.service.ts (điểm thưởng và thành tích) & 30 & 10 \\ \hline
    Toàn dự án (ngưỡng tối thiểu chung) & 28 & 14 \\ \hline
\end{longtable}

Mỗi lần thực thi npm test hoặc npm run test:report, một custom Jest reporter đặt tại test/reporters/markdown-reporter.js tự động sinh tài liệu docs/testing/test-results.md, ghi nhận đầy đủ ID, mô tả, trạng thái và thời gian thực thi của từng test case. Tài liệu này là nguồn dữ liệu tham chiếu cho các bảng kết quả ở các mục tiếp theo.

\subsection{Phạm vi kiểm thử}

Hai mươi sáu file spec của backend-core được tổ chức thành 16 nhóm phù hợp với cấu trúc mô-đun của hệ thống. Mười lăm nhóm đầu thuộc tầng unit và integration (25 file, tổng 207 test case); nhóm thứ mười sáu là bộ E2E gồm 12 test case mô phỏng toàn bộ luồng xác thực trên cơ sở dữ liệu PostgreSQL thật. Bảng~\ref{tab:test_scope} liệt kê chi tiết. Cột File spec hiển thị tên file rút gọn; đường dẫn đầy đủ với specs mô-đun là backend-core/src/modules/<module>/tests/, còn các bộ test dùng chung nằm dưới backend-core/test/.

% 5 cột p{}, divisor=100: 4+18+32+6+40; -4pt buffer cho rounding
\setlength{\tbu}{\dimexpr(\linewidth - 10\tabcolsep - 6\arrayrulewidth - 4pt)/100\relax}
\begin{longtable}{|p{4\tbu}|p{18\tbu}|p{32\tbu}|p{6\tbu}|p{40\tbu}|}
    \caption{Phạm vi kiểm thử backend-core}
    \label{tab:test_scope} \\
    \hline
    \textbf{STT} & \textbf{Nhóm chức năng} & \textbf{File spec} & \textbf{Số TC} & \textbf{Trọng tâm kiểm thử} \\ \hline
    \endfirsthead
    \multicolumn{5}{r}{\textit{(tiếp trang trước)}} \\
    \hline
    \textbf{STT} & \textbf{Nhóm chức năng} & \textbf{File spec} & \textbf{Số TC} & \textbf{Trọng tâm kiểm thử} \\ \hline
    \endhead
    \hline \multicolumn{5}{r}{\textit{(tiếp trang sau)}} \\
    \endfoot
    \hline
    \endlastfoot

    1 & Xác thực và phân quyền (auth) & register.spec.ts\newline login.spec.ts\newline auth-extended.spec.ts\newline jwt-strategy.spec.ts & 29 & Đăng ký với DTO validation và bcrypt hash; đăng nhập + cấp JWT; refresh token; đổi mật khẩu; Google OAuth ID token; JwtStrategy của Passport \\ \hline

    2 & Cập nhật thông tin tài khoản (users) & update-profile.spec.ts & 12 & DTO whitelist (forbidNonWhitelisted), email uniqueness, role enum, trạng thái hoạt động \\ \hline

    3 & Cộng đồng -- đăng bài (posts) & posts.spec.ts & 11 & CRUD bài đăng, like/unlike, ẩn bài đã xóa, kích hoạt gamification onEvent (XP = 5) \\ \hline

    4 & Thông báo (notifications) & notifications.spec.ts & 9 & Phân trang, ownership filter, đánh dấu tất cả đã đọc, xóa thông báo \\ \hline

    5 & Luyện phát âm IPA (pronunciation) & pronunciation.spec.ts\newline pronunciation.service.spec.ts & 21 & CRUD âm vị (RolesGuard chỉ ADMIN), Redis cache và invalidation, đánh giá mức độ thành thạo IPA, vòng đời attempt từ PENDING đến hoàn tất \\ \hline

    6 & IELTS -- thống kê, advanced, chuỗi học (ielts) & ielts-statistics.spec.ts\newline ielts-advanced.service.spec.ts\newline streak.service.spec.ts & 40 & Thống kê theo từng giai đoạn (Foundation/Basic/Advanced/Intensive), vòng đời L/R/W/S session, đánh giá đáp án, dispatch tác vụ AI qua RabbitMQ, chuỗi học liên tục và mốc thành tích \\ \hline

    7 & Luyện nghe chép chính tả (dictation) & dictation-progress.service.spec.ts & 6 & Tiến độ upsert, sentence XP, lesson completion notification, không award trùng \\ \hline

    8 & Luyện shadowing video (shadowing) & shadowing-progress.service.spec.ts & 6 & Tiến độ upsert, sentence XP, hoàn thành bài khi đủ số câu \\ \hline

    9 & Vocab Lab -- lặp lại ngắt quãng (vocab-lab) & vocab-lab-fsrs.service.spec.ts & 10 & Thuật toán FSRS với các mức rating (Good/Again/Easy), lịch ôn theo trường due, ownership deck \\ \hline

    10 & Gamification -- XP và thành tích (gamification) & gamification.service.spec.ts & 10 & Tạo XpLog, cập nhật totalXp, thông báo lên level, grant achievement, không tạo trùng \\ \hline

    11 & Bài thi thử IELTS Intensive (exams) & exams-submit.spec.ts\newline result-callback.spec.ts & 10 & Submit + dispatch publishGradingTask, quota guard cho gói FREE, HMAC signature callback từ AI worker, idempotency \\ \hline

    12 & Polling kết quả chấm AI (results) & results.spec.ts & 6 & Trạng thái GRADING/GRADED/GRADING\_FAILED; bảo vệ IDOR (trả 404 ẩn existence) \\ \hline

    13 & Subscription và thanh toán (subscriptions) & subscriptions-webhook.spec.ts\newline subscriptions.service.spec.ts\newline subscriptions.cron.spec.ts & 13 & VNPay IPN HMAC-SHA512, replay protection, dùng thử 7 ngày, cancel subscription, quota usage và rollback, downgrade cron sau grace period 3 ngày \\ \hline

    14 & Edge case -- phân trang và giá trị biên (test/edge-cases) & pagination.spec.ts\newline boundary-values.spec.ts & 16 & Query param bất thường (page = 0/-1, limit cực lớn, cursor không tồn tại), điểm 0/100/101, body rỗng, mảng imageUrls quá kích thước \\ \hline

    15 & Bảo mật -- truy cập chéo người dùng (test/security) & cross-user-access.spec.ts & 8 & IDOR trên VocabLab/Notes/Posts (404 vs 403 tùy mô hình dữ liệu), khe hở admin bypass đã được ghi nhận \\ \hline

    16 & E2E xác thực -- toàn bộ chu trình (test/e2e) & auth.e2e-spec.ts & 12 & Postgres thật trên schema test, đủ luồng register/login/refresh; kèm các kịch bản token hết hạn, token bị giả mạo chữ ký, đăng ký trùng email \\ \hline
\end{longtable}

Hai nhóm Edge case (STT~14) và Security (STT~15) không gắn với một mô-đun đơn lẻ mà kiểm thử xuyên suốt nhiều endpoint, tập trung vào các loại lỗi thường gặp như IDOR (Insecure Direct Object Reference), kiểm tra giá trị biên và phân trang bất thường. Nhóm Subscription (STT~13) bao gồm cả webhook IPN của cổng thanh toán VNPay (một luồng tích hợp với hệ thống bên ngoài) nên được tách thành ba file spec để phân tách trách nhiệm: kiểm thử webhook, lifecycle service và cron downgrade.

\subsection{Danh sách các test case}

Trên cơ sở phạm vi đã xác định, toàn bộ 207 test case unit/integration của backend-core được trình bày chi tiết trong bảng~\ref{tab:testcase_list}. Bảng có sáu cột: mã định danh, số thứ tự trường hợp, chức năng kiểm tra, tiền điều kiện, tình huống thực hiện và kết quả mong muốn. Các trường hợp kế thừa từ phiên bản trước (TC01 đến TC11) đại diện cho bảy chức năng cốt lõi -- Đăng ký, Đăng nhập, Đăng bài, Thông báo, Cập nhật thông tin tài khoản, Luyện phát âm và Thống kê IELTS -- được mô tả theo góc nhìn người dùng cuối. Các test case bổ sung sử dụng quy ước đặt mã có dạng TC\_TIỀN-TỐ-MÔ-ĐUN\_SỐ-TRƯỜNG-HỢP và mô tả ở mức kỹ thuật để phản ánh chính xác hành vi mà từng test thực hiện, bao gồm các luồng xác thực mở rộng (TC\_AUTH, TC\_JWT), PronunciationService (TC\_PRO\_S), IELTS Advanced service (TC\_IADV), thuật toán lặp lại ngắt quãng FSRS (TC\_FSRS), gamification (TC\_GAMI), chuỗi học (TC\_STREAK), tiến độ dictation và shadowing (TC\_DCT, TC\_SHD), bài thi IELTS Intensive (TC\_EXAM, TC\_CB, TC\_RES), subscription thanh toán (TC\_SUB), giá trị biên (TC\_BOUND), phân trang bất thường (TC\_EDGE) và bảo mật truy cập chéo (TC\_SEC).

% 6 cột p{}, divisor=134: (1.0+0.4+2.0+2.8+3.5+3.7)×10; -4pt buffer cho rounding
\setlength{\tbu}{\dimexpr(\linewidth - 12\tabcolsep - 7\arrayrulewidth - 4pt)/134\relax}
\begin{longtable}{|p{10\tbu}|p{4\tbu}|p{20\tbu}|p{28\tbu}|p{35\tbu}|p{37\tbu}|}
    \caption{Danh sách test case}
    \label{tab:testcase_list} \\
    \hline
    \textbf{ID} & \textbf{TH} & \textbf{Chức năng} & \textbf{Tiền điều kiện} & \textbf{Tình huống test} & \textbf{Kết quả mong muốn} \\ \hline
    \endfirsthead
    \multicolumn{6}{r}{\textit{(tiếp trang trước)}} \\
    \hline
    \textbf{ID} & \textbf{TH} & \textbf{Chức năng} & \textbf{Tiền điều kiện} & \textbf{Tình huống test} & \textbf{Kết quả mong muốn} \\ \hline
    \endhead
    \hline \multicolumn{6}{r}{\textit{(tiếp trang sau)}} \\
    \endfoot
    \hline
    \endlastfoot

    TC01 & 1 & Đăng ký & Chưa có tài khoản, ở trang Đăng ký & Để trống họ tên người dùng & Thông báo họ tên không được bỏ trống \\ \hline
    TC01 & 2 & Đăng ký & Chưa có tài khoản, ở trang Đăng ký & Nhập email sai định dạng & Thông báo email không hợp lệ \\ \hline
    TC01 & 3 & Đăng ký & Chưa có tài khoản, ở trang Đăng ký & Nhập mật khẩu ít hơn 6 ký tự & Thông báo mật khẩu phải có ít nhất 6 ký tự \\ \hline
    TC01 & 4 & Đăng ký & Email đã đăng ký trước đó & Nhập email đã tồn tại trong hệ thống & Thông báo email đã được sử dụng \\ \hline
    TC01 & 5 & Đăng ký & Chưa có tài khoản, ở trang Đăng ký & Nhập đầy đủ thông tin hợp lệ & Tạo tài khoản thành công, trả thông tin tài khoản (không hiển thị mật khẩu) \\ \hline
    TC01 & 6 & Đăng ký & Chưa có tài khoản, ở trang Đăng ký & Nhập mật khẩu hợp lệ và nộp đăng ký & Mật khẩu được lưu ở dạng mã hoá, không lưu nguyên văn \\ \hline
    TC01 & 7 & Đăng ký & Chưa có tài khoản, ở trang Đăng ký & Đăng ký thành công với thông tin hợp lệ & Hệ thống tự động tạo bộ thẻ từ vựng ``Default'' cho tài khoản mới \\ \hline
    TC01 & 8 & Đăng ký & Chưa có tài khoản, ở trang Đăng ký & Đăng ký không chọn vai trò / chọn vai trò hợp lệ (Học viên hoặc Giáo viên) & Gán vai trò Học viên mặc định / lưu đúng vai trò được chọn \\ \hline
    TC01 & 9 & Đăng ký & Chưa có tài khoản, ở trang Đăng ký & Đăng ký với vai trò Quản trị viên & Từ chối yêu cầu, thông báo vai trò không hợp lệ \\ \hline

    TC02 & 1 & Đăng nhập & Tài khoản tồn tại, ở trang Đăng nhập & Email đúng, mật khẩu sai & Thông báo thông tin đăng nhập không chính xác \\ \hline
    TC02 & 2 & Đăng nhập & Tài khoản tồn tại, ở trang Đăng nhập & Email không tồn tại, mật khẩu đúng & Thông báo thông tin đăng nhập không chính xác \\ \hline
    TC02 & 3 & Đăng nhập & Tài khoản tồn tại, ở trang Đăng nhập & Email và mật khẩu đều sai & Thông báo thông tin đăng nhập không chính xác \\ \hline
    TC02 & 4 & Đăng nhập & Tài khoản tồn tại, ở trang Đăng nhập & Chỉ nhập email, để trống mật khẩu & Thông báo thông tin đăng nhập không chính xác \\ \hline
    TC02 & 5 & Đăng nhập & Tài khoản tồn tại, ở trang Đăng nhập & Email và mật khẩu đúng & Đăng nhập thành công, hệ thống cấp mã truy cập kèm thông tin tài khoản \\ \hline

    TC03 & 1 & Đăng bài & Đã đăng nhập, ở Newsfeed & Đăng bài với nội dung để trống & Thông báo cần nhập nội dung \\ \hline
    TC03 & 2 & Đăng bài & Đã đăng nhập, ở Newsfeed & Đăng bài với nội dung vượt 10000 ký tự & Thông báo nội dung vượt giới hạn cho phép \\ \hline
    TC03 & 3 & Đăng bài & Đã đăng nhập, ở Newsfeed & Chọn loại bài viết không thuộc danh mục cho phép & Thông báo loại bài viết không hợp lệ \\ \hline
    TC03 & 4 & Đăng bài & Chưa đăng nhập & Cố gắng đăng bài với nội dung hợp lệ & Yêu cầu đăng nhập, không cho phép đăng bài \\ \hline
    TC03 & 5 & Đăng bài & Đã đăng nhập, ở Newsfeed & Nhập nội dung hợp lệ và nộp bài & Đăng bài thành công, hệ thống cộng điểm kinh nghiệm cho người đăng \\ \hline
    TC03 & 6 & Đăng bài & Đã đăng nhập, ở Newsfeed & Đăng bài kèm danh sách thẻ và hình ảnh & Bài đăng được lưu đầy đủ thẻ và hình ảnh đã đính kèm \\ \hline
    TC03 & 7 & Xem bài & Đã đăng nhập, ở Newsfeed & Mở danh sách bài viết & Hiển thị bài viết theo thứ tự mới nhất, ẩn các bài đã bị gỡ \\ \hline
    TC03 & 8 & Tương tác & Đã đăng nhập, bài chưa được thích & Nhấn Thích lên bài viết & Bài được đánh dấu đã thích, số lượt thích tăng thêm 1 \\ \hline
    TC03 & 9 & Tương tác & Đã đăng nhập, bài đã thích trước đó & Nhấn Thích lần thứ hai lên cùng bài viết & Bỏ thích, số lượt thích giảm đi 1 \\ \hline
    TC03 & 10 & Tương tác & Đã đăng nhập & Nhấn Thích lên bài viết đã bị xoá & Thông báo không tìm thấy bài viết \\ \hline
    TC03 & 11 & Xem bài & Đã đăng nhập & Truyền tham số phân trang sai định dạng & Hệ thống từ chối yêu cầu phân trang \\ \hline

    TC05 & 1 & Thông báo & Chưa đăng nhập & Truy cập danh sách thông báo & Yêu cầu đăng nhập \\ \hline
    TC05 & 2 & Thông báo & Chưa đăng nhập & Truy cập số lượng thông báo chưa đọc & Yêu cầu đăng nhập \\ \hline
    TC05 & 3 & Thông báo & Chưa đăng nhập & Xoá một thông báo & Yêu cầu đăng nhập, không thực hiện xoá \\ \hline
    TC05 & 4 & Thông báo & Đã đăng nhập & Mở danh sách thông báo trang đầu & Hiển thị 20 thông báo mới nhất của tài khoản, sắp xếp giảm dần \\ \hline
    TC05 & 5 & Thông báo & Đã đăng nhập & Mở danh sách thông báo trang thứ hai & Hiển thị 20 thông báo tiếp theo, đúng quy tắc phân trang \\ \hline
    TC05 & 6 & Thông báo & Đã đăng nhập & Đánh dấu một thông báo của mình là đã đọc & Thông báo được cập nhật trạng thái đã đọc \\ \hline
    TC05 & 7 & Thông báo & Đã đăng nhập & Cố gắng đánh dấu thông báo của người dùng khác & Hệ thống bỏ qua, không cập nhật và không lộ thông tin \\ \hline
    TC05 & 8 & Thông báo & Đã đăng nhập & Nhấn Đánh dấu tất cả là đã đọc & Toàn bộ thông báo chưa đọc của tài khoản chuyển sang đã đọc \\ \hline
    TC05 & 9 & Thông báo & Đã đăng nhập & Xoá một thông báo của chính mình & Thông báo bị xoá khỏi danh sách \\ \hline

    TC06 & 1 & Cập nhật thông tin & Chưa đăng nhập & Cập nhật thông tin cá nhân & Yêu cầu đăng nhập \\ \hline
    TC06 & 2 & Cập nhật thông tin & Đã đăng nhập, ở trang cá nhân & Nhập email mới sai định dạng & Thông báo email không hợp lệ \\ \hline
    TC06 & 3 & Cập nhật thông tin & Đã đăng nhập, ở trang cá nhân & Chọn vai trò ngoài ba vai trò được hỗ trợ & Thông báo vai trò không hợp lệ \\ \hline
    TC06 & 4 & Cập nhật thông tin & Đã đăng nhập, ở trang cá nhân & Nhập trạng thái hoạt động không phải đúng/sai & Thông báo dữ liệu không hợp lệ \\ \hline
    TC06 & 5 & Cập nhật thông tin & Đã đăng nhập, ở trang cá nhân & Nhập họ tên dưới dạng số & Thông báo dữ liệu không hợp lệ \\ \hline
    TC06 & 6 & Cập nhật thông tin & Đã đăng nhập, ở trang cá nhân & Gửi kèm trường dữ liệu không được hỗ trợ & Thông báo trường dữ liệu không hợp lệ \\ \hline
    TC06 & 7 & Cập nhật thông tin & Đã đăng nhập, ở trang cá nhân & Cập nhật họ tên hợp lệ & Cập nhật thành công, thông tin mới được lưu lại \\ \hline
    TC06 & 8 & Cập nhật thông tin & Đã đăng nhập, ở trang cá nhân & Cập nhật email mới chưa được sử dụng & Cập nhật thành công, email mới được lưu \\ \hline
    TC06 & 9 & Cập nhật thông tin & Đã đăng nhập, ở trang cá nhân & Cập nhật email đã có tài khoản khác sử dụng & Thông báo email đã được sử dụng \\ \hline
    TC06 & 10 & Cập nhật thông tin & Đã đăng nhập, ở trang cá nhân & Cập nhật vai trò sang Giáo viên & Cập nhật thành công, vai trò mới được lưu \\ \hline
    TC06 & 11 & Cập nhật thông tin & Đã đăng nhập, ở trang cá nhân & Tắt trạng thái hoạt động của tài khoản & Cập nhật thành công, tài khoản chuyển sang ngừng hoạt động \\ \hline
    TC06 & 12 & Cập nhật thông tin & Đã đăng nhập, ở trang cá nhân & Gửi yêu cầu cập nhật với dữ liệu rỗng & Hệ thống không thay đổi thông tin, giữ nguyên trạng thái hiện tại \\ \hline

    TC10 & 1 & Luyện phát âm & Không yêu cầu đăng nhập & Mở danh sách âm vị công khai & Hiển thị danh sách âm vị, người chưa đăng nhập vẫn xem được \\ \hline
    TC10 & 2 & Luyện phát âm & Không yêu cầu đăng nhập & Tra cứu một âm vị không tồn tại & Thông báo không tìm thấy âm vị \\ \hline
    TC10 & 3 & Luyện phát âm & Không yêu cầu đăng nhập & Tra cứu một âm vị có trong cơ sở dữ liệu & Hiển thị đầy đủ thông tin âm vị \\ \hline
    TC10 & 4 & Luyện phát âm & Chưa đăng nhập & Gửi kết quả luyện phát âm lên hệ thống & Yêu cầu đăng nhập, không ghi nhận kết quả \\ \hline
    TC10 & 5 & Luyện phát âm & Đã đăng nhập, hoàn thành một bài luyện & Gửi điểm phát âm hợp lệ lên hệ thống & Hệ thống ghi nhận tiến độ luyện phát âm cho tài khoản \\ \hline
    TC10 & 6 & Quản lý âm vị & Đã đăng nhập với vai trò Học viên & Tạo mới một âm vị & Từ chối truy cập, yêu cầu vai trò Quản trị viên \\ \hline
    TC10 & 7 & Quản lý âm vị & Đã đăng nhập với vai trò Quản trị viên & Tạo âm vị mới với đầy đủ thông tin bắt buộc & Tạo âm vị thành công \\ \hline
    TC10 & 8 & Quản lý âm vị & Đã đăng nhập với vai trò Quản trị viên & Tạo âm vị nhưng thiếu thông tin bắt buộc & Thông báo dữ liệu không hợp lệ \\ \hline

    TC11 & 1 & Thống kê & Chưa đăng nhập & Xem thống kê tổng quan & Yêu cầu đăng nhập \\ \hline
    TC11 & 2 & Thống kê & Chưa đăng nhập & Xem thống kê giai đoạn Foundation & Yêu cầu đăng nhập \\ \hline
    TC11 & 3 & Thống kê & Đã đăng nhập & Xem thống kê tổng quan tiến độ học & Hiển thị thống kê tổng quan của tài khoản \\ \hline
    TC11 & 4 & Thống kê & Đã đăng nhập & Xem thống kê giai đoạn Foundation & Hiển thị thống kê giai đoạn Foundation của tài khoản \\ \hline
    TC11 & 5 & Thống kê & Đã đăng nhập & Xem thống kê giai đoạn Basic & Hiển thị thống kê giai đoạn Basic của tài khoản \\ \hline
    TC11 & 6 & Thống kê & Đã đăng nhập & Xem thống kê giai đoạn Advanced & Hiển thị thống kê giai đoạn Advanced của tài khoản \\ \hline
    TC11 & 7 & Thống kê & Đã đăng nhập & Xem thống kê giai đoạn Intensive & Hiển thị thống kê giai đoạn Intensive của tài khoản \\ \hline

    % ===== 146 test case bổ sung =====

    % Nhóm TC_AUTH — refresh token, đổi mật khẩu, Google OAuth (10 case)
    TC\_AUTH\_03 & 1 & Refresh token & Đã đăng nhập, có refresh token hợp lệ & Gửi POST /auth/refresh kèm refresh token hợp lệ & 201, trả về cặp accessToken và refreshToken mới \\ \hline
    TC\_AUTH\_03 & 2 & Refresh token & Refresh token đã hết hạn & Gửi POST /auth/refresh với refresh token expired & 401 Unauthorized \\ \hline
    TC\_AUTH\_03 & 3 & Refresh token & Refresh token bị giả mạo chữ ký & Gửi POST /auth/refresh với refresh token tampered & 401 Unauthorized \\ \hline
    TC\_AUTH\_04 & 1 & Đổi mật khẩu & Đã đăng nhập, nhập đúng currentPassword & Gửi POST /auth/change-password với mật khẩu hiện tại chính xác & 201 kèm thông báo ``changed successfully'', mật khẩu mới được hash bcrypt \\ \hline
    TC\_AUTH\_04 & 2 & Đổi mật khẩu & Đã đăng nhập, nhập sai currentPassword & Gửi POST /auth/change-password với mật khẩu hiện tại sai & 400 ``Current password is incorrect'' \\ \hline
    TC\_AUTH\_04 & 3 & Đổi mật khẩu & Tài khoản đăng ký qua Google OAuth (password = null) & Gửi POST /auth/change-password trên tài khoản Google & 400, từ chối vì tài khoản đăng nhập bằng Google \\ \hline
    TC\_AUTH\_05 & 1 & Google OAuth & Chưa có tài khoản trên hệ thống & Gửi POST /auth/google với Google ID token hợp lệ (user mới) & 201, tạo user mới kèm bộ thẻ Default cho Vocab Lab \\ \hline
    TC\_AUTH\_05 & 2 & Google OAuth & Đã có user với googleId khớp & Gửi POST /auth/google với Google ID token của user cũ & 201, đăng nhập bình thường, không tạo user mới \\ \hline
    TC\_AUTH\_05 & 3 & Google OAuth & Google ID token có audience không thuộc danh sách cho phép & Gửi POST /auth/google với token sai audience & 401 ``Invalid Google ID token'' \\ \hline
    TC\_AUTH\_05 & 4 & Google OAuth & Google ID token hợp lệ nhưng payload thiếu email & Gửi POST /auth/google với payload không có email & 400 Bad Request \\ \hline

    % Nhóm TC_JWT — JwtStrategy của Passport (5 case)
    TC\_JWT & 1 & JwtStrategy validate & Payload JWT đầy đủ sub, email, role & Passport gọi validate() với payload hợp lệ & Trả về object \{id, email, role\} cho request scope \\ \hline
    TC\_JWT & 2 & JwtStrategy validate & Token được ký bằng secret khác JWT\_SECRET & Gửi request kèm token sai secret & Passport từ chối, ném UnauthorizedException \\ \hline
    TC\_JWT & 3 & JwtStrategy validate & Token có trường exp ở thời điểm quá khứ & Gửi request kèm token expired & Passport từ chối, ném UnauthorizedException \\ \hline
    TC\_JWT & 4 & JwtStrategy validate & Payload JWT thiếu trường sub & validate() được gọi với payload không có sub & Trả id = undefined (khe hở đã ghi nhận, cần ném UnauthorizedException) \\ \hline
    TC\_JWT & 5 & JwtStrategy validate & Payload JWT thiếu trường role & validate() được gọi với payload không có role & Trả role = undefined (khe hở đã ghi nhận, cần ném UnauthorizedException) \\ \hline

    % Nhóm TC_PRO_S — PronunciationService (13 case)
    TC\_PRO\_S & 1 & PronunciationService -- đọc danh sách & Redis cache đã có danh sách sounds & Gọi getAllSounds() & Trả kết quả từ cache, không truy vấn Prisma \\ \hline
    TC\_PRO\_S & 2 & PronunciationService -- đọc danh sách & Redis cache rỗng & Gọi getAllSounds() & Truy vấn DB qua Prisma và ghi kết quả vào cache \\ \hline
    TC\_PRO\_S & 3 & PronunciationService -- đọc theo symbol & Cache có sound theo symbol cụ thể & Gọi getSoundBySymbol(symbol) & Trả từ cache, không gọi Prisma \\ \hline
    TC\_PRO\_S & 4 & PronunciationService -- đọc theo symbol & Cache miss & Gọi getSoundBySymbol(symbol) & Truy vấn DB và ghi vào cache \\ \hline
    TC\_PRO\_S & 5 & PronunciationService -- thêm sound & Admin gọi createSound với payload hợp lệ & Service lưu sound mới vào DB & Sound được tạo, cache bị invalidate \\ \hline
    TC\_PRO\_S & 6 & PronunciationService -- cập nhật sound & Admin gọi updateSound & Service cập nhật trường của sound & Sound được update, cache bị invalidate \\ \hline
    TC\_PRO\_S & 7 & PronunciationService -- xóa sound & Admin gọi deleteSound & Service xóa sound khỏi DB & Sound bị xóa, cache bị invalidate \\ \hline
    TC\_PRO\_S & 8 & PronunciationService -- progress user & User có progress trên một số sound & Gọi getUserProgress(userId) & Trả progress đã join với toàn bộ danh sách sounds \\ \hline
    TC\_PRO\_S & 9 & PronunciationService -- thống kê & User có lịch sử attempt phong phú & Gọi getUserStats(userId) & Tính các chỉ số mastery cho từng loại âm \\ \hline
    TC\_PRO\_S & 10 & PronunciationService -- progress theo từ & User có attempt cho các từ ví dụ & Gọi getWordProgress(userId) & Trả danh sách từ kèm mastery status dựa trên attempt \\ \hline
    TC\_PRO\_S & 11 & PronunciationService -- cập nhật progress & User submit điểm phát âm hợp lệ & Gọi updateProgress(userId, soundId, score) & Upsert progress, invalidate cache, gọi gamification.onEvent để cộng XP \\ \hline
    TC\_PRO\_S & 12 & PronunciationService -- tạo attempt & User bắt đầu một lượt luyện phát âm & Gọi createPronunciationAttempt & Tạo bản ghi attempt với trạng thái PENDING \\ \hline
    TC\_PRO\_S & 13 & PronunciationService -- cập nhật attempt & Attempt PENDING đã tồn tại & Gọi updatePronunciationAttempt với chi tiết & Cập nhật các trường của attempt \\ \hline

    % Nhóm TC_FSRS — Vocab Lab spaced repetition (10 case)
    TC\_FSRS\_01 & 1 & VocabLab -- submitReview rating Good & User có card mới chưa từng review & Gọi submitReview với rating = Good (3) & Card được lưu với stability lớn hơn 0 và due lớn hơn now \\ \hline
    TC\_FSRS\_01 & 2 & VocabLab -- submitReview rating Good & Như trường hợp trên & submitReview hoàn tất & Tạo một bản ghi FlashcardReview gắn với card \\ \hline
    TC\_FSRS\_02 & 1 & VocabLab -- submitReview rating Again & User có card đã review trước đó & Gọi submitReview với rating = Again (1) & Stability reset về xấp xỉ 0.4 (mặc định FSRS) \\ \hline
    TC\_FSRS\_02 & 2 & VocabLab -- submitReview rating Again & Card mới & submitReview với Again & due được lập lịch trong khoảng 10 phút sau \\ \hline
    TC\_FSRS\_03 & 1 & VocabLab -- so sánh interval & Cùng một card được test với hai rating khác nhau & submitReview với rating Easy (4) so với Good (3) & scheduledDays khi Easy luôn lớn hơn hoặc bằng scheduledDays khi Good \\ \hline
    TC\_FSRS\_04 & 1 & VocabLab -- ownership card & Card thuộc deck của user khác & submitReview với userId không khớp deck.userId & Throw NotFoundException \\ \hline
    TC\_FSRS\_04 & 2 & VocabLab -- card không tồn tại & cardId không có trong DB & submitReview với cardId sai & Throw NotFoundException \\ \hline
    TC\_FSRS\_05 & 1 & VocabLab -- getStudyCards & Deck có cards NEW + LEARNING/REVIEW/RELEARNING với due đã đến hạn & Gọi getStudyCards(deckId) & Trả về cards NEW cùng cards đến hạn ôn \\ \hline
    TC\_FSRS\_05 & 2 & VocabLab -- getStudyCards loại trừ tương lai & Có cards với due lớn hơn now & Gọi getStudyCards(deckId) & Loại trừ những cards có lịch ôn tương lai \\ \hline
    TC\_FSRS\_05 & 3 & VocabLab -- getStudyCards ownership & Deck không tồn tại hoặc không thuộc user & Gọi getStudyCards(deckId, userId) & Throw NotFoundException \\ \hline

    % Nhóm TC_GAMI — gamification XP & achievement (10 case)
    TC\_GAMI\_01 & 1 & Gamification onEvent (xp dương) & Event gửi tới có xp dương & Service xử lý event & Tạo XpLog với đầy đủ trường (userId, eventKey, xp, reason) \\ \hline
    TC\_GAMI\_01 & 2 & Gamification onEvent (xp dương) & Như trên, IeltsProfile của user tồn tại & Service xử lý event & Cập nhật totalXp của IeltsProfile \\ \hline
    TC\_GAMI\_01 & 3 & Gamification onEvent (xp dương) & IeltsProfile của user chưa tồn tại & Service xử lý event & Không crash, log nhẹ và bỏ qua bước cập nhật profile \\ \hline
    TC\_GAMI\_02 & 1 & Gamification level-up & XP mới vượt ngưỡng nâng level & Service xử lý event & Gửi notification loại ACHIEVEMENT cho user \\ \hline
    TC\_GAMI\_02 & 2 & Gamification level-up & XP tăng nhưng cùng mức level & Service xử lý event & KHÔNG gửi notification level-up \\ \hline
    TC\_GAMI\_03 & 1 & Gamification achievement trùng & User đã có achievement trước đó & Service grant lại cùng achievement key & Bỏ qua UserAchievement.create, không tạo bản ghi trùng \\ \hline
    TC\_GAMI\_04 & 1 & Gamification onEvent (xp = 0) & Event không kèm XP & Service xử lý event & Bỏ qua tạo XpLog hoàn toàn \\ \hline
    TC\_GAMI\_04 & 2 & Gamification onEvent (xp = 0) & Event có achievementKeys nhưng xp = 0 & Service xử lý event & Vẫn xử lý achievementKeys bình thường \\ \hline
    TC\_GAMI\_05 & 1 & Gamification tryGrantAchievement & User chưa có achievement & Gọi tryGrantAchievement(userId, key) & Tạo UserAchievement và gửi notification ACHIEVEMENT \\ \hline
    TC\_GAMI\_05 & 2 & Gamification tryGrantAchievement & achievementKey không có định nghĩa trong DB & Gọi tryGrantAchievement với key sai & No-op, không tạo gì cả \\ \hline

    % Nhóm TC_STREAK — chuỗi học liên tục (13 case)
    TC\_STREAK\_01 & 1 & Streak recordActivity (lần đầu) & IeltsProfile chưa có lastActiveDate & Gọi recordActivity(userId) & currentStreak = 1, lastActiveDate = ngày hôm nay \\ \hline
    TC\_STREAK\_01 & 2 & Streak recordActivity (lần đầu) & Như trường hợp trên & Gọi recordActivity(userId) & longestStreak ít nhất bằng 1 \\ \hline
    TC\_STREAK\_02 & 1 & Streak recordActivity (liên tiếp) & lastActiveDate = ngày hôm qua & Gọi recordActivity(userId) & currentStreak tăng thêm 1 \\ \hline
    TC\_STREAK\_02 & 2 & Streak recordActivity (liên tiếp) & Streak mới vượt longestStreak cũ & Gọi recordActivity(userId) & Cập nhật longestStreak bằng giá trị mới \\ \hline
    TC\_STREAK\_03 & 1 & Streak recordActivity (gap 2+ ngày) & lastActiveDate cách hôm nay từ 2 ngày trở lên & Gọi recordActivity(userId) & currentStreak reset về 1 \\ \hline
    TC\_STREAK\_03 & 2 & Streak recordActivity (gap 2+ ngày) & Như trên & Gọi recordActivity(userId) & longestStreak KHÔNG bị giảm \\ \hline
    TC\_STREAK\_04 & 1 & Streak recordActivity (cùng ngày) & lastActiveDate = hôm nay & Gọi recordActivity(userId) & Trả về profile, KHÔNG gọi update (idempotent) \\ \hline
    TC\_STREAK\_04 & 2 & Streak recordActivity (cùng ngày) & Như trên & Gọi recordActivity(userId) & KHÔNG kích hoạt gamification event \\ \hline
    TC\_STREAK\_05 & 1 & Streak milestone 7 ngày & Streak mới đạt 7 ngày liên tiếp & Gọi recordActivity (lần thứ 7) & Gọi GamificationService.onEvent kèm achievementKeys chứa XM\_ON\_FIRE \\ \hline
    TC\_STREAK\_05 & 2 & Streak milestone 7 ngày & Như trên & Gọi recordActivity & Gửi milestone notification (STREAK\_MILESTONES có 7) \\ \hline
    TC\_STREAK\_05 & 3 & Streak milestone 7 ngày & Như trên & Gọi recordActivity & XP truyền vào tỷ lệ với streak mới (5 nhân newStreak) \\ \hline
    TC\_STREAK\_06 & 1 & Streak getStreakStats edge & IeltsProfile không tồn tại & Gọi getStreakStats(userId) & Trả về null \\ \hline
    TC\_STREAK\_06 & 2 & Streak getStreakStats edge & Prisma throw error & Gọi getStreakStats(userId) & Trả null, KHÔNG throw (silent fail) \\ \hline

    % Nhóm TC_DCT — DictationProgressService (6 case)
    TC\_DCT & 1 & DictationProgress findByLesson & Progress chưa có trong DB & Gọi findByLesson(userId, lessonId) & Trả completedSentences = [] và difficulty mặc định \\ \hline
    TC\_DCT & 2 & DictationProgress findByLesson & Progress đã có trong DB & Gọi findByLesson(userId, lessonId) & Trả completedSentences và difficulty từ DB \\ \hline
    TC\_DCT & 3 & DictationProgress upsert & Progress mới, số câu hoàn thành nhỏ hơn total & Gọi upsert(userId, lessonId, completed) & Tạo bản ghi, award sentence XP, KHÔNG complete lesson \\ \hline
    TC\_DCT & 4 & DictationProgress upsert & Đủ số câu (bằng total) & Gọi upsert & Chuyển sang complete, gửi notification, award lesson achievement \\ \hline
    TC\_DCT & 5 & DictationProgress upsert & Lesson đã hoàn thành trước đó & Gọi upsert lại & KHÔNG kích hoạt completion notification lần thứ hai \\ \hline
    TC\_DCT & 6 & DictationProgress findAllByUser & User có progress trên nhiều lesson & Gọi findAllByUser(userId) & Trả map progress theo lessonId \\ \hline

    % Nhóm TC_SHD — ShadowingProgressService (6 case)
    TC\_SHD & 1 & ShadowingProgress findByLesson & Progress chưa có trong DB & Gọi findByLesson(userId, lessonId) & Trả completedSentences = [] \\ \hline
    TC\_SHD & 2 & ShadowingProgress findByLesson & Progress đã có trong DB & Gọi findByLesson(userId, lessonId) & Trả completedSentences từ DB \\ \hline
    TC\_SHD & 3 & ShadowingProgress upsert & Số câu nhỏ hơn tổng số câu video & Gọi upsert & Award sentence XP, KHÔNG award lesson complete \\ \hline
    TC\_SHD & 4 & ShadowingProgress upsert & Đủ tổng số câu của video & Gọi upsert & Complete video, award lesson completion achievement \\ \hline
    TC\_SHD & 5 & ShadowingProgress upsert & newCount không lớn hơn existingCount & Gọi upsert & KHÔNG kích hoạt sentence XP (no-op) \\ \hline
    TC\_SHD & 6 & ShadowingProgress findAllByUser & User có progress trên nhiều video & Gọi findAllByUser(userId) & Trả map progress theo lessonId \\ \hline

    % Nhóm TC_IADV — IeltsAdvancedService (20 case)
    TC\_IADV & 1 & IELTS Advanced -- danh sách Listening & Có nhiều Listening parts trong DB & Gọi getListeningParts() & Danh sách trả về sắp xếp theo partNumber \\ \hline
    TC\_IADV & 2 & IELTS Advanced -- chi tiết Listening part & partId không tồn tại & Gọi getListeningPartDetail(partId) & Throw NotFoundException \\ \hline
    TC\_IADV & 3 & IELTS Advanced -- submit Listening & User submit đáp án nhiều định dạng (multiple choice, fill-in, matching) & Gọi submitListeningPart & Đánh giá đúng từng định dạng, tạo session lưu kết quả \\ \hline
    TC\_IADV & 4 & IELTS Advanced -- high-score XP Listening & Tổng điểm Listening đạt từ 80\% trở lên & Gọi submitListeningPart & Thưởng thêm XP cho high-score \\ \hline
    TC\_IADV & 5 & IELTS Advanced -- danh sách Reading & Có nhiều Reading parts trong DB & Gọi getReadingParts() & Sắp xếp theo partNumber \\ \hline
    TC\_IADV & 6 & IELTS Advanced -- chi tiết Reading part & partId không tồn tại & Gọi getReadingPartDetail(partId) & Throw NotFoundException \\ \hline
    TC\_IADV & 7 & IELTS Advanced -- submit Reading & User submit đáp án Reading & Gọi submitReadingPart & Đánh giá đúng, award reading achievements \\ \hline
    TC\_IADV & 8 & IELTS Advanced -- danh sách Writing prompts & Có nhiều writing prompts trong DB & Gọi getWritingPrompts (phân trang) & Trả danh sách prompts kèm best score mapping \\ \hline
    TC\_IADV & 9 & IELTS Advanced -- chi tiết Writing prompt & User có active session và lịch sử bài viết & Gọi getWritingPromptDetail(promptId) & Trả chi tiết prompt cùng active session và history \\ \hline
    TC\_IADV & 10 & IELTS Advanced -- tạo Writing session & User đã có session IN\_PROGRESS & Gọi createWritingSession & Trả session đang có, không tạo session mới \\ \hline
    TC\_IADV & 11 & IELTS Advanced -- tạo Writing session & User chưa có session đang chạy & Gọi createWritingSession & Tạo session mới \\ \hline
    TC\_IADV & 12 & IELTS Advanced -- lưu draft Writing & Session đang active & Gọi saveWritingDraft & Cập nhật draft của session đang active \\ \hline
    TC\_IADV & 13 & IELTS Advanced -- lưu draft Writing & Session không active hoặc không tồn tại & Gọi saveWritingDraft & Throw NotFoundException \\ \hline
    TC\_IADV & 14 & IELTS Advanced -- submit Writing & User submit bài Writing & Gọi submitWritingSession & Status chuyển sang GRADING, publish task lên RabbitMQ cho AI worker \\ \hline
    TC\_IADV & 15 & IELTS Advanced -- tạo Speaking session & User chưa hoặc đang có session & Gọi createSpeakingSession & Tạo mới hoặc trả session đang chạy \\ \hline
    TC\_IADV & 16 & IELTS Advanced -- submit Speaking & Không có audio answer nào được upload & Gọi submitSpeakingSession & Fail / từ chối submit \\ \hline
    TC\_IADV & 17 & IELTS Advanced -- submit Speaking & Có đầy đủ audio answers hợp lệ & Gọi submitSpeakingSession & Lưu answers, publish task AI cho RabbitMQ \\ \hline
    TC\_IADV & 18 & IELTS Advanced -- thống kê & User có nhiều session đã hoàn thành & Gọi getStatistics(userId) & Tổng hợp điểm đúng qua các session \\ \hline
    TC\_IADV & 19 & IELTS Advanced -- lịch sử chi tiết & sessionId thuộc user yêu cầu & Gọi getHistoryDetail(sessionId, userId) & Trả chi tiết session Listening \\ \hline
    TC\_IADV & 20 & IELTS Advanced -- lịch sử chi tiết & sessionId thuộc user khác & Gọi getHistoryDetail(sessionId, otherUserId) & Throw NotFoundException (bảo vệ ownership) \\ \hline

    % Nhóm TC_EXAM — Submit bài thi IELTS Intensive (6 case)
    TC\_EXAM & 1 & Submit bài thi IELTS Intensive & User có session IN\_PROGRESS (WRITING), còn quota & POST /exams/sessions/:id/submit hợp lệ & 201, status = SUBMITTED, publishGradingTask được gọi đúng 1 lần kèm đúng payload \\ \hline
    TC\_EXAM & 2 & Submit bài thi IELTS Intensive & Session đã ở trạng thái GRADED & POST submit lại session đã chấm xong & 400 ``Session already graded'' \\ \hline
    TC\_EXAM & 3 & Submit bài thi IELTS Intensive & sessionId thuộc user khác & POST submit với JWT của user khác & 403 Access denied \\ \hline
    TC\_EXAM & 4 & Submit bài thi IELTS Intensive & Quota AI\_WRITING\_GRADING vượt giới hạn gói FREE & POST submit khi đã hết quota tháng & 403 QUOTA\_EXCEEDED \\ \hline
    TC\_EXAM & 5 & Submit bài thi IELTS Intensive & Không có JWT token & POST submit không có header Authorization & 401 Unauthorized \\ \hline
    TC\_EXAM & 6 & Submit bài thi IELTS Intensive & sessionId không tồn tại & POST submit với sessionId sai & 404 Not Found \\ \hline

    % Nhóm TC_CB — Result callback từ AI worker (4 case)
    TC\_CB & 1 & Callback chấm AI & AI worker hoàn tất grading, gửi callback với HMAC signature đúng & POST callback với signature hợp lệ & 200 OK, lưu IeltsIntensiveResult, session.status = GRADED \\ \hline
    TC\_CB & 2 & Callback chấm AI & Signature không khớp HMAC & POST callback với signature giả mạo & 401 Unauthorized, KHÔNG ghi vào DB \\ \hline
    TC\_CB & 3 & Callback chấm AI & Session đã ở trạng thái GRADING\_FAILED & POST callback với cùng sessionId & 200 OK (idempotent), result không bị ghi đè \\ \hline
    TC\_CB & 4 & Callback chấm AI & Body thiếu totalScore và feedback & POST callback với payload không đủ trường & 400 Bad Request \\ \hline

    % Nhóm TC_RES — Polling kết quả chấm AI (6 case)
    TC\_RES & 1 & Polling kết quả chấm & Session đang ở GRADING & GET /results/:sessionId & 200, trả \{status: ``GRADING'', result: null\} \\ \hline
    TC\_RES & 2 & Polling kết quả chấm & Session đã GRADED & GET /results/:sessionId & 200, trả full result shape \\ \hline
    TC\_RES & 3 & Polling kết quả chấm & sessionId thuộc user khác & GET /results/:sessionId với JWT user khác & 404 (không tiết lộ existence) \\ \hline
    TC\_RES & 4 & Polling kết quả chấm & sessionId không tồn tại & GET /results/:sessionId với id sai & 404 Not Found \\ \hline
    TC\_RES & 5 & Polling kết quả chấm & Không có JWT token & GET /results/:sessionId không kèm token & 401 Unauthorized \\ \hline
    TC\_RES & 6 & Polling kết quả chấm & Session ở GRADING\_FAILED & GET /results/:sessionId & 200, trả \{status: ``GRADING\_FAILED'', error: string\} \\ \hline

    % Nhóm TC_SUB — Subscription, VNPay IPN và cron downgrade (13 case)
    TC\_SUB & 1 & VNPay IPN -- checksum sai & vnp\_SecureHash sai (HMAC không khớp) & GET /subscriptions/webhook/vnpay với hash sai & RspCode ``97'' (Fail checksum) \\ \hline
    TC\_SUB & 2 & VNPay IPN -- thanh toán thành công & ResponseCode = 00, HMAC đúng & GET webhook payment thành công & payment status = ``succeeded'', subscription nâng lên PREMIUM \\ \hline
    TC\_SUB & 3 & VNPay IPN -- user hủy giao dịch & vnp\_ResponseCode = 24 (user hủy) & GET webhook cancel & RspCode ``00'' non-success, payment.create KHÔNG được gọi \\ \hline
    TC\_SUB & 4 & VNPay IPN -- replay & Session đã xóa, payment đã tồn tại & GET webhook lại với cùng TxnRef & RspCode ``02'', KHÔNG double-credit \\ \hline
    TC\_SUB & 5 & VNPay IPN -- thiếu TxnRef & Payload thiếu vnp\_TxnRef & GET webhook & RspCode ``99'' (Missing txnRef) \\ \hline
    TC\_SUB & 6 & SubscriptionsService -- startTrial & User chưa từng dùng trial & Gọi startTrial(userId) & Tạo subscription TRIALING, endDate xấp xỉ now cộng 7 ngày \\ \hline
    TC\_SUB & 7 & SubscriptionsService -- startTrial & trialUsed = true & Gọi startTrial(userId) & Throw BadRequestException ``already used'' \\ \hline
    TC\_SUB & 8 & SubscriptionsService -- cancel & User đang ở gói FREE & Gọi cancelSubscription(userId) & Throw BadRequestException \\ \hline
    TC\_SUB & 9 & SubscriptionsService -- cancel & User đang ở gói PREMIUM & Gọi cancelSubscription(userId) & status = CANCELED, gọi paymentProvider.cancelSubscription \\ \hline
    TC\_SUB & 10 & SubscriptionsService -- incrementUsage & Còn quota AI\_WRITING\_GRADING & Gọi incrementUsage(AI\_WRITING\_GRADING) & Trả về true, count tăng thêm 1 \\ \hline
    TC\_SUB & 11 & SubscriptionsService -- incrementUsage & Đã vượt quota & Gọi incrementUsage(AI\_WRITING\_GRADING) & Trả về false, count bị rollback \\ \hline
    TC\_SUB & 12 & Subscriptions Cron -- downgrade & Subscription quá endDate cộng 3 ngày grace period & Cron 2:00 AM chạy & Downgrade về FREE \\ \hline
    TC\_SUB & 13 & Subscriptions Cron -- downgrade & Trong grace period (endDate nhỏ hơn now, dưới 3 ngày) & Cron chạy & KHÔNG downgrade, giữ trạng thái \\ \hline

    % Nhóm TC_EDGE — Edge case pagination (8 case)
    TC\_EDGE & 1 & Notifications pagination & User đã đăng nhập & GET /notifications?page=0 & 200 (ParseInt không reject 0) \\ \hline
    TC\_EDGE & 2 & Notifications pagination & Như trên & GET /notifications?page=-1 & 200 (ParseInt chấp nhận số âm, không có @Min) \\ \hline
    TC\_EDGE & 3 & Notifications pagination & Như trên & GET /notifications?limit=10000 & 200 (không có cap maxLimit) \\ \hline
    TC\_EDGE & 4 & Notifications pagination & Như trên & GET /notifications?cursor=not-a-valid-uuid & 200 (notifications không có cursor param, bị whitelist drop) \\ \hline
    TC\_EDGE & 5 & Notifications pagination & Như trên & GET /notifications không có query param & 200, DefaultValuePipe gán page=1, limit=20 \\ \hline
    TC\_EDGE & 6 & Notifications pagination & Như trên & GET /notifications?limit=0 & 200 (ParseInt chấp nhận 0, không có @Min(1)) \\ \hline
    TC\_EDGE & 7 & Notifications empty & User không có notification nào trong DB & GET /notifications & 200, trả \{notifications: [], total: 0\}, KHÔNG throw 404 \\ \hline
    TC\_EDGE & 8 & Posts cursor pagination & cursor đúng định dạng UUID nhưng không tồn tại & GET /posts?cursor=uuid-không-tồn-tại & 200, trả \{items: []\} \\ \hline

    % Nhóm TC_BOUND — Boundary value (8 case)
    TC\_BOUND & 1 & Pronunciation score boundary & User đăng nhập & POST /pronunciation/progress với score = 0 & 201 (lower boundary, @Min(0) pass) \\ \hline
    TC\_BOUND & 2 & Pronunciation score boundary & User đăng nhập & POST với score = 100 & 201 (upper boundary hợp lệ) \\ \hline
    TC\_BOUND & 3 & Pronunciation score boundary & User đăng nhập & POST với score = 101 & 201 (over max, hiện tại pass vì DTO thiếu @Max(100), khe hở đã ghi nhận) \\ \hline
    TC\_BOUND & 4 & Pronunciation score boundary & User đăng nhập & POST với score = -1 & 400 (under @Min(0), ValidationPipe reject) \\ \hline
    TC\_BOUND & 5 & IELTS Statistics zero-data & User chưa học gì & GET /ielts-statistics/foundation & 200, trả \{vocabulary: \{wordsLearned: 0\}\}, KHÔNG throw \\ \hline
    TC\_BOUND & 6 & Posts body boundary & User đăng nhập & POST /posts với body = chuỗi rỗng & 400 (@MinLength(1) violation) \\ \hline
    TC\_BOUND & 7 & Posts imageUrls boundary & User đăng nhập & POST /posts với imageUrls = [] & 201 (mảng rỗng được phép) \\ \hline
    TC\_BOUND & 8 & Posts imageUrls boundary & User đăng nhập & POST /posts với imageUrls có 11 phần tử & 201 (hiện tại pass vì thiếu @ArrayMaxSize(10), khe hở đã ghi nhận) \\ \hline

    % Nhóm TC_SEC — Cross-user access control (8 case)
    TC\_SEC & 1 & Cross-user VocabLab & userB đăng nhập, deck thuộc userA & PATCH /vocab-lab/decks/:id (của userA) với token userB & 404 (ẩn existence theo mô hình ownership) \\ \hline
    TC\_SEC & 2 & Cross-user VocabLab & userB đăng nhập, deck thuộc userA & DELETE /vocab-lab/decks/:id với token userB & 404 \\ \hline
    TC\_SEC & 3 & Cross-user VocabLab & userB đăng nhập, deck thuộc userA & GET /vocab-lab/decks/:id/cards với token userB & 404 \\ \hline
    TC\_SEC & 4 & Cross-user Notes & userB đăng nhập, note thuộc userA & DELETE /notes/:id với token userB & 404 \\ \hline
    TC\_SEC & 5 & Cross-user Notes & userB đăng nhập, note thuộc userA & PATCH /notes/:id với token userB & 404 \\ \hline
    TC\_SEC & 6 & Cross-user Posts & userB (non-admin) đăng nhập, post thuộc userA & DELETE /posts/:id với token userB & 403 (Posts có phân quyền role rõ ràng nên trả 403 thay vì 404) \\ \hline
    TC\_SEC & 7 & Cross-user Posts & userB đăng nhập, post thuộc userA & PATCH /posts/:id với token userB & 404 \\ \hline
    TC\_SEC & 8 & Admin bypass Posts & admin (role ADMIN) cố xóa post của user khác & DELETE /posts/:id với token admin & Mong đợi 200 (admin được phép). Hiện tại trả 403, khe hở admin bypass đã ghi nhận \\ \hline

\end{longtable}

\subsection{Kết quả kiểm thử}

Sau khi hoàn thiện hệ thống, toàn bộ 207 test case unit/integration cùng 12 test case E2E được thực thi tự động qua hai lệnh: npm run test:report (cho unit/integration) và npm run test:e2e (cho E2E). Cả hai lệnh đều kích hoạt Jest và ghi log chi tiết per-test vào tài liệu auto-generated docs/testing/test-results.md. Lần chạy gần nhất ngày 20/05/2026 cho thấy tất cả 207 test case unit/integration đều đạt trạng thái \textbf{Pass} sau 63,7 giây thực thi; 12 test case E2E (chạy trên PostgreSQL schema test) cũng đều đạt \textbf{Pass} hoàn toàn. Bảng~\ref{tab:testcase_result} tổng hợp kết quả theo từng nhóm chức năng đã định nghĩa ở phạm vi kiểm thử; chi tiết per-test có thể đối chiếu trong tài liệu docs/testing/test-results.md đính kèm với mã nguồn.

\renewcommand{\arraystretch}{1.2}
% 8 cột p{}, divisor=100: 4+28+8+8+8+10+20+14; -4pt buffer cho rounding
\setlength{\tbu}{\dimexpr(\linewidth - 16\tabcolsep - 9\arrayrulewidth - 4pt)/100\relax}
\begin{longtable}{|p{4\tbu}|p{28\tbu}|p{8\tbu}|p{8\tbu}|p{8\tbu}|p{10\tbu}|p{20\tbu}|p{14\tbu}|}
    \caption{Tổng hợp kết quả kiểm thử theo nhóm}
    \label{tab:testcase_result} \\
    \hline
    \textbf{STT} & \textbf{Nhóm chức năng} & \textbf{Số TC} & \textbf{Pass} & \textbf{Fail} & \textbf{Tỷ lệ} & \textbf{Người TH} & \textbf{Ngày} \\ \hline
    \endfirsthead
    \multicolumn{8}{r}{\textit{(tiếp trang trước)}} \\
    \hline
    \textbf{STT} & \textbf{Nhóm chức năng} & \textbf{Số TC} & \textbf{Pass} & \textbf{Fail} & \textbf{Tỷ lệ} & \textbf{Người TH} & \textbf{Ngày} \\ \hline
    \endhead
    \hline \multicolumn{8}{r}{\textit{(tiếp trang sau)}} \\
    \endfoot
    \hline
    \endlastfoot

    1 & Xác thực và phân quyền & 29 & 29 & 0 & 100\% & Auto (Jest) & 20/05/2026 \\ \hline
    2 & Cập nhật thông tin tài khoản & 12 & 12 & 0 & 100\% & Auto (Jest) & 20/05/2026 \\ \hline
    3 & Cộng đồng -- đăng bài & 11 & 11 & 0 & 100\% & Auto (Jest) & 20/05/2026 \\ \hline
    4 & Thông báo & 9 & 9 & 0 & 100\% & Auto (Jest) & 20/05/2026 \\ \hline
    5 & Luyện phát âm IPA & 21 & 21 & 0 & 100\% & Auto (Jest) & 20/05/2026 \\ \hline
    6 & IELTS (thống kê, advanced, chuỗi học) & 40 & 40 & 0 & 100\% & Auto (Jest) & 20/05/2026 \\ \hline
    7 & Luyện nghe chép chính tả & 6 & 6 & 0 & 100\% & Auto (Jest) & 20/05/2026 \\ \hline
    8 & Luyện shadowing video & 6 & 6 & 0 & 100\% & Auto (Jest) & 20/05/2026 \\ \hline
    9 & Vocab Lab (lặp lại ngắt quãng) & 10 & 10 & 0 & 100\% & Auto (Jest) & 20/05/2026 \\ \hline
    10 & Gamification (XP và thành tích) & 10 & 10 & 0 & 100\% & Auto (Jest) & 20/05/2026 \\ \hline
    11 & Bài thi thử IELTS Intensive & 10 & 10 & 0 & 100\% & Auto (Jest) & 20/05/2026 \\ \hline
    12 & Polling kết quả chấm AI & 6 & 6 & 0 & 100\% & Auto (Jest) & 20/05/2026 \\ \hline
    13 & Subscription và thanh toán & 13 & 13 & 0 & 100\% & Auto (Jest) & 20/05/2026 \\ \hline
    14 & Edge case (phân trang, giá trị biên) & 16 & 16 & 0 & 100\% & Auto (Jest) & 20/05/2026 \\ \hline
    15 & Bảo mật (truy cập chéo người dùng) & 8 & 8 & 0 & 100\% & Auto (Jest) & 20/05/2026 \\ \hline
    16 & E2E xác thực (toàn bộ chu trình) & 12 & 12 & 0 & 100\% & Auto (Jest) & 20/05/2026 \\ \hline
\end{longtable}

Kết quả này khẳng định backend-core hiện đáp ứng đầy đủ các yêu cầu nghiệp vụ đã đặc tả ở chương~3, đồng thời đảm bảo các ràng buộc bảo mật (IDOR, JWT lifecycle), tính nhất quán cấp cơ sở dữ liệu (unique email, ownership deck/note/post) và tính chính xác của các luồng tích hợp ngoài như VNPay IPN và AI grading callback. Năm khe hở nhỏ được phát hiện trong quá trình kiểm thử (TC\_JWT\_04, TC\_JWT\_05, TC\_BOUND\_03, TC\_BOUND\_08, TC\_SEC\_08) là các trường hợp test hiện đạt PASS theo hành vi mặc định nhưng đã ghi nhận điểm cần cải thiện cho phiên bản kế tiếp -- chi tiết được nêu trong cột ``Kết quả mong muốn'' của các test case tương ứng ở bảng~\ref{tab:testcase_list}.

